\section{Overview of the proof}\label{sec:overview}

This section describes the argument in words, fixes the terminology and the
parameters, and works through the simulation of one inference. It contains no
proofs; every claim made here is proved in the section cited for it.

Fix a \(\Res\) refutation of \(\BPHP^{n+1}_n\) with \(S\) nodes. The proof
has three steps.

\subsection{Step 1: from a refutation to low-degree polynomial calculus}
A linear clause \(C\) is a disjunction of affine equations. Write each
equation through its \emph{true-indicator} \(g_i\), the affine form that is
one exactly when the equation holds, so that \(C\) is falsified exactly on
the affine flat \(Z(C)=\{g_1=\cdots=g_s=0\}\). Following Buss, Impagliazzo,
Kraj\'i\v{c}ek, Pudl\'ak, Razborov, and Sgall \cite{BIKPRS}, we attach to
\(C\) fresh Boolean variables \(r_{ui}\) and the product
\[
 P_C=\prod_{u=1}^{h}\Bigl(1-\sum_{i=1}^{s}r_{ui}g_i\Bigr),
\]
together with the axioms \(g_iP_C=0\) and \(r_{ui}^2=r_{ui}\). On \(Z(C)\)
the product is one. At a point where \(C\) holds, each factor vanishes for
half of the choices of its coefficients, so \(P_C\) behaves like the
indicator of ``\(C\) is false'' up to an error that halves with each of the
\(h\) factors. The axioms \(g_iP_C=0\) say that \(P_C\) is supported on
\(Z(C)\). The line \(P_C=0\) is \emph{not} an axiom: deriving it is how the
simulation expresses that \(C\) has been proved.

What matters is that \(P_C\) has degree \(2h\) whatever the width or the
rank of \(C\). The identity \(1-P_C=\sum_iU_{C,i}g_i\)
(Corollary~\ref{lem:prefix}) lets a few local polynomial identities of
degree at most \(4h+1\) simulate complementary-parity resolution and
semantic weakening (Lemmas~\ref{lem:clause-resolution} and
\ref{lem:clause-weakening}). Each simulated inference uses the \emph{lines}
\(P_A=0\), \(P_B=0\) derived earlier, never their derivations, so the
degree does not grow along the proof. This is why neither the depth of
the refutation nor regularity enters: the simulation is local and the
proof DAG is only followed in topological order. An initial clause costs
degree \(2h+\ell\). Altogether (Theorem~\ref{thm:clause-transfer}) the
refutation becomes a polynomial calculus refutation of degree
\[
 D=\max\{2h+\ell,\,4h+1\}
\]
from the compact pigeonhole axioms \(\mathcal Q_{m,\ell}\) of
\eqref{eq:compact-base} together with at most
\(N_{\mathrm{reg}}=3S+\binom m2\) blocks of extension axioms. With
\(h=3\ell\) this is \(D=12\ell+1=O(\log n)\).

\subsection{Step 2: removing all extension variables at once}
A degree-\(O(\log n)\) refutation with extension variables contradicts
nothing by itself. We substitute polynomials of degree at most \(k\) in the
original bit variables for all the \(r_{ui}\), simultaneously, and replay
the refutation (Lemma~\ref{lem:substitution}). A block whose inputs have
rank \(r\le h(k+1)\) can be substituted \emph{exactly}: its \(h\) factors,
each with coefficients of degree at most \(k\), suffice to turn \(P_C\) into
the true indicator of \(Z(C)\), and the extension axioms become consequences
of the Boolean axioms (Lemma~\ref{lem:packing}). This is the origin of the
rank threshold \(h(k+1)\).

A block of higher rank cannot be treated this way within degree \(k\).
Instead we look for a polynomial \(f\) of degree at most \(k\) whose
restriction to the flat \(Z(C)\) is identically zero. Then
\(f=\sum_ia_ig_i\) with \(\deg a_i<k\), and substituting the \(a_i\) into
the first factor and zero into the others turns \(P_C\) into \(1-f\). The
extension axiom \(g_iP_C\) becomes \(g_i(1-f)\), which is not a consequence
of anything; but after multiplication by \(f\) it becomes
\(-g_i(f^2-f)\), a consequence of the Boolean axioms. So we replay the
whole refutation \emph{multiplied by \(f\)}. Every axiom image is then
derivable from \(\mathcal Q_{m,\ell}\) alone, and the final line \(1\)
becomes \(f\) (Lemma~\ref{lem:hybrid}):
\[
 f\in\Cons_{B}(\mathcal Q_{m,\ell}),\qquad B=k(D+1).
\]
This is why the end product is a derivation of a nonzero polynomial \(f\)
and not of the constant one: the weight \(f\) has to vanish on the
falsifying flats of all high-rank clauses at once, and the constant one
does not.

A single \(f\) must serve every high-rank block. A flat of codimension
\(r>h(k+1)\) has so few free coordinates that the restrictions of
degree-\(k\) polynomials to it span a space smaller, by a factor
\(e^{-k^2/m}\), than the space \(\mathcal L_k\) of \emph{row-linear}
polynomials (at most one bit variable from each pigeon). If
\[
 m\ln(4M)\le k^2,
\]
where \(M\) bounds the number of high-rank blocks, a dimension count gives a
nonzero \(f\in\mathcal L_k\) restricting to zero on all of them
(Lemma~\ref{lem:kernel}). This inequality is the only place where the size
of the refutation enters, through \(M\le N_{\mathrm{reg}}=3S+\binom m2\).

\subsection{Step 3: a nonzero row-linear polynomial is not derivable}
It remains to show that no nonzero \(f\in\mathcal L_k\) has a polynomial
calculus derivation of degree \(B\) from the pigeonhole axioms when \(B\)
is below roughly \(n/2\) (Theorem~\ref{thm:row-separation}). This is
stronger than the degree lower bound for refutations, which is the case
\(f=1\). We decode bit labels into the unary (one variable per pigeon and
hole) functional pigeonhole principle, where, as in Razborov's analysis
\cite{Raz98}, the linear functionals vanishing on all degree-\(B\)
consequences are exactly the solutions of the matching-moment equations
\eqref{eq:marginals}. Theorem~\ref{thm:moments} shows that every solution
through a lower degree extends through degree \(B\) when \(2B-1\) is at
most the number of holes; the extension problem is a cycle-filling problem
in chessboard complexes, solved by the vanishing of their homology in the
range of Bj\"orner, Lov\'asz, Vre\'cica, and \v{Z}ivaljevi\'c
\cite{BLVZ94}. To separate a given \(f\) of degree \(t\), a signed sum over a
\(t\)-dimensional cube of partial assignments isolates one top-degree
coefficient of \(f\) as the constant one, while mapping pigeonhole
consequences to consequences of a smaller pigeonhole instance of degree
\(B-t\), where a normalized functional exists.

\subsection{Parameters}
The three steps meet in the inequalities of Theorem~\ref{thm:affine}:
\[
 (n+1)\ln(4M)\le k^2,\qquad 2k(D+1)-1\le n,\qquad 4(k-1)<n .
\]
With \(D=12\ell+1\), the second allows \(k\) up to about \(n/(24\ell)\),
and then the first holds as long as \(\ln S\) is below a constant times
\(n/\ell^2\). This gives Theorem~\ref{thm:main}. The symbols used
throughout are collected here.
\begin{center}
\begin{tabular}{@{}p{0.27\textwidth}p{0.68\textwidth}@{}}
\toprule
\(\ell,\ n=2^\ell,\ m=n+1,\ v=m\ell\) & bits per label, holes, pigeons (rows), bit variables\\
\(N\) & number of columns of a unary system (\(N=n\) in the application)\\
\(S,\ N_{\mathrm{reg}}=3S+\binom m2\) & nodes of the refutation; blocks (slots) of the registry\\
\(h\) & number of factors of a block product (its \emph{accuracy}); \(h=3\ell\)\\
\(s\) & number of inputs of a block (clause width after compression, at most \(v+2\))\\
\(r\) & rank of a block: dimension of the span of its inputs; threshold \(h(k+1)\)\\
\(D\) & degree of the simulating PC refutation; \(D=12\ell+1\)\\
\(k\) & degree bound for \(f\) and for the substituted coefficients\\
\(t\) & degree of a particular \(f\); number of rows in the separating cube; \(t\le k\)\\
\(M\) & bound on the number of proper high-rank blocks; \(M\le N_{\mathrm{reg}}\)\\
\(B=k(D+1)\) & degree at which \(f\) would be derived; needs \(2B-1\le n\)\\
\bottomrule
\end{tabular}
\end{center}

\subsection{Terminology}
Some terms come from the research notebook in which the proof was developed
and are kept because the formal statements use them.
\emph{Old} variables are the bit variables \(b_{it}\) of the formula (in
Section~\ref{sec:moments}, the unary variables), as opposed to the fresh
coefficient variables \(r_{ui}\); the \emph{old base} is the axiom system
before blocks are added, and an \emph{old-base consequence} is a polynomial
derivable from it alone. A \emph{complete affine extension block of
accuracy \(h\)} is the package \eqref{eq:block}: affine inputs \(g_i\) in
the old variables, fresh coefficients, all \emph{companions} \(g_iP\), and
all coefficient Boolean equations; ``complete'' means that every companion
is present, and ``one-level'' that no block's inputs mention another
block's variables. A \emph{registry} is a family of blocks fixed in advance,
one \emph{slot} per clause that the simulation may need. A polynomial is
\emph{row-linear} if each monomial has at most one bit variable from each
pigeon. \(\Ideal_B\) and \(\Cons_B\) denote degree-\(B\) Nullstellensatz and
polynomial calculus consequences (Section~\ref{sec:algebra}); ``NS'' and
``PC'' abbreviate the two systems.

\subsection{A worked example: one resolution step}
Take accuracy \(h=1\) and resolve \(A=(u=1)\) with \(B=(u=0)\), for an
affine form \(u\), to the empty clause. The true-indicators are \(u\) for
\(A\) and \(1-u\) for \(B\), so with coefficients \(r\) and \(s\)
\[
 P_A=1-ru,\qquad P_B=1-s(1-u),\qquad P_\varnothing=1 .
\]
Suppose the lines \(P_A\) and \(P_B\) have been derived. Since
\(u^2-u\) is a combination of Boolean axioms, polynomial calculus derives
\[
 (1-u)P_A-r(u^2-u)=(1-u)(1-ru)-r(u^2-u)=1-u ,
\]
which is \eqref{eq:pivot-correction} in this case: from ``\(A\) holds'' we
have derived that the indicator of \(B\) is zero. Multiplying by \(s\) and
adding \(P_B\) gives
\[
 s(1-u)+\bigl(1-s(1-u)\bigr)=1 ,
\]
the value of the empty clause, and the refutation is complete. All lines
have degree at most three. With side literals present, the same
computation is carried out after multiplying by the product \(P_T\) of the
conclusion \(T\), and the side literals are absorbed by the companions of
\(T\); with \(h\) factors, \(r\) and \(s\) are replaced by the polynomials
\(U_{A,u}\) and \(U_{B,1-u}\) of degree \(2h-1\), which gives the bound
\(4h+1\) of Lemma~\ref{lem:clause-resolution}.

\subsection{Why the restriction-based approaches stop short}
As we understand the earlier lower bounds for regular and bounded-depth
\(\Res\) \cite{EGI25,BCD24,BC25,BI25,BCBI26,AI25,EI25}, they follow a path
through the refutation and maintain a set of linear forms (a closure, or
the state of a game or of a random walk) that records what the path has
learned. Regularity or a depth bound controls how much such a path can
accumulate; a general DAG may reuse a node along many long paths, and the
control is lost. The present argument never follows a path. It is static
and algebraic: the whole DAG is translated into one polynomial calculus
refutation whose degree does not depend on depth, in the manner of the
extension-variable method of \cite{BIKPRS}. That method was designed for
bounded-depth Frege systems with counting gates, where it reduces lower
bounds to degree bounds for systems with extension axioms; in that
generality no low-degree way of removing the extension variables is known. For
\(\Res\) the extension axioms have inputs that are affine in the original
variables, and this is exactly what makes the one-shot removal of Step~2
possible: low rank allows exact substitution, and high rank leaves room
for a common annihilating polynomial. The price is that the base system
must exclude every nonzero row-linear consequence, not only the constant
one, which is the Razborov-style statement of Step~3. For nested
extension axioms, as they arise from \(\mathrm{AC}^0[p]\)-Frege proofs, we
know of no corresponding removal step.
